\documentclass[10pt]{article}
\usepackage[T1]{fontenc}
\usepackage[utf8]{inputenc}
\usepackage[french]{babel}
\usepackage{amsmath,amssymb,mathrsfs,mathpazo}
\usepackage{geometry}
\geometry{a4paper,margin=2cm}
\usepackage{graphicx,array,tabularx,tabularray,colortbl,multirow,multicol,booktabs,enumitem,pifont,ulem,fancybox,diagbox,cancel}
\usepackage{pstricks,pstricks-add,pst-plot,pst-tree,pst-node,pst-text,pst-func,pst-math,pst-bspline,pst-3dplot,pst-eucl,pst-solides3d}
\let\pscalculate\undefined
\usepackage{pst-bezier}
\usepackage{tikz,pgfplots,tkz-tab}
\usetikzlibrary{arrows,arrows.meta,patterns,calc,positioning,trees,decorations.markings,shapes.geometric,shapes.misc,shapes.arrows,intersections,decorations.pathmorphing,decorations.pathreplacing,backgrounds,shadows,babel}
\pgfplotsset{compat=1.18}
\usepgfplotslibrary{fillbetween}
\usepackage[np]{numprint}
\usepackage{listings,makecell,pas-tableur,twoopt}
\let\barre\undefined
\usepackage[pstricks]{tablvar}
\definecolor{darkgreen}{RGB}{0,100,0}
\providecommand{\red}{\color{red}}
\providecommand{\blue}{\color{blue}}
\providecommand{\green}{\color{green}}
\providecommand{\black}{\color{black}}
\providecommand{\white}{\color{white}}
\providecommand{\gray}{\color{gray}}
\providecommand{\cyan}{\color{cyan}}
\providecommand{\magenta}{\color{magenta}}
\providecommand{\yellow}{\color{yellow}}
\providecommand{\violet}{\color{violet}}
\providecommand{\orange}{\color{orange}}
\providecommand{\euro}{\texteuro}
\providecommand{\EUR}{\texteuro}
\providecommand{\ieme}{\textsuperscript{e}}
\providecommand{\degres}{\ensuremath{^\circ}}
\providecommand{\degree}{\ensuremath{^\circ}}
\providecommand{\vv}[1]{\overrightarrow{#1}}
\providecommand{\Coord}[1]{#1}
\providecommand{\decosix}{}
\providecommand{\starredbullet}{\ensuremath{\bullet}}
\pagestyle{empty}
\setlength{\parindent}{0pt}
\usepackage{adjustbox}\DeclareUnicodeCharacter{2236}{:}\DeclareUnicodeCharacter{2190}{\ensuremath{\leftarrow}}
\providecommand{\R}{\mathbb{R}}
\providecommand{\N}{\mathbb{N}}
\providecommand{\Z}{\mathbb{Z}}
\providecommand{\Q}{\mathbb{Q}}
\providecommand{\D}{\mathbb{D}}
\providecommand{\C}{\mathbb{C}}
\providecommand{\e}{{\mathrm{e}}}
\providecommand{\vect}[1]{\overrightarrow{#1}}
\providecommand{\vectt}[1]{\overrightarrow{\mathrm{#1}}}
\providecommand{\vvt}[1]{\overrightarrow{\mathrm{#1}}}
\providecommand{\np}[2][]{{#2~#1}}
\providecommand{\vv}[1]{{\overrightarrow{#1}}}
\renewcommand{\d}{{\mathrm{d}}}
\providecommand{\degres}{{^\circ}}
\providecommand{\degre}{{^\circ}}
\providecommand{\blue}{\color{blue}}
\providecommand{\red}{\color{red}}
\providecommand{\enskip}{\quad}
\providecommand{\llbracket}{\lbrack\!\lbrack}
\providecommand{\rrbracket}{\rbrack\!\rbrack}
\providecommand{\euro}{{\text{€}}}
\providecommand{\no}{{\text{n°}}}
\providecommand{\interval}[2]{{[#1;#2]}}
\providecommand{\intervFO}[2]{{[#1;#2[}}
\providecommand{\dotfill}{\ldots}

\providecommand{\ds}{\displaystyle}
\providecommand{\textding}[1]{#1}
\providecommand{\Oij}{\(\mathrm{O};\vec{\imath},\vec{\jmath}\)}
\providecommand{\Oijk}{\(\mathrm{O};\vec{\imath},\vec{\jmath},\vec{k}\)}
\providecommand{\Ouv}{\(\mathrm{O};\vec{u},\vec{v}\)}
\providecommand{\intervFF}[2]{\ensuremath{[#1;#2]}}
\renewcommand{\euro}{\ifmmode\text{\texteuro}\else\texteuro\fi}
\DeclareUnicodeCharacter{00B0}{\ifmmode^{\circ}\else\textdegree\fi}
\DeclareUnicodeCharacter{20AC}{\ifmmode\text{\texteuro}\else\texteuro\fi}
% Énoncés : Éducation nationale. Transcription et corrigés : APMEP, Denis Vergès et contributeurs. Découpage et classement : CoopMaths, Sébastien Lozano, Éric Elter et contributeurs. Adaptation de la présentation : Maths Explorer.
% CC BY-SA : https://coopmaths.fr/www/about/
% https://coopmaths.fr/alea/static/bac/2021/tex/bac_2021_09_metropole_9.tex
\begin{document}
\noindent\textit{APMEP / CoopMaths. CC BY-SA. Présentation adaptée par Maths Explorer.}\par\medskip
\medskip

\begin{tabular}{|l|}\hline
Principaux domaines abordés :\\
Suites numériques ; raisonnement par récurrence.\\ \hline
\end{tabular}

\medskip

On considère les suites $\left(u_n\right)$ et $\left(v_n\right)$ définies par:

\[u_0 = 16 \quad ;\quad  v_0 = 5 \;;\]

et pour tout entier naturel $n$ :

\renewcommand\arraystretch{2}
\[\left\{\begin{array}{l c l}
u_{n+1}&=&\dfrac{3u_n + 2v_n}{5}\\
v_{n+1}&=& \dfrac{u_n + v_n}{2}
\end{array}\right.\]
\renewcommand\arraystretch{1}

\smallskip

\begin{enumerate}
\item Calculer $u_1$ et $v_1$.
\item On considère la suite $\left(w_n\right)$ définie pour tout entier naturel $n$ par : $w_n = u_n - v_n$.
	\begin{enumerate}
		\item Démontrer que la suite $\left(w_n\right)$ est géométrique de raison 0,1.
		
En déduire, pour tout entier naturel $n$, l'expression de $w_n$ en fonction de $n$.
		\item Préciser le signe de la suite $\left(w_n\right)$  et la limite de cette suite.
	\end{enumerate}
\item 
	\begin{enumerate}
		\item Démontrer que, pour tout entier naturel $n$, on a : $u_{n+1} - u_n = - 0,4 w_n$.
		\item En déduire que la suite $\left(u_n\right)$ est décroissante.

\end{enumerate}
		
On peut démontrer de la même manière que la suite $\left(v_n\right)$ est croissante. On admet ce résultat, et on remarque qu'on a alors: pour tout entier naturel $n$,\, $v_n \geqslant v_0 = 5$.
\begin{enumerate}[resume]
		\item Démontrer par récurrence que, pour tout entier naturel $n$, on a : $u_n \geqslant 5$. 
		
En déduire que la suite $\left(u_n\right)$ est convergente. On appelle $\ell$ la limite de $\left(u_n\right)$.
	\end{enumerate}
\end{enumerate}
	
On peut démontrer de la même manière que la suite $\left(v_n\right)$ est convergente. On admet ce résultat, et on appelle $\ell'$ la limite de $\left(v_n\right)$.
\begin{enumerate}[resume]
\item 
	\begin{enumerate}
		\item Démontrer que $\ell = \ell'$.
		\item On considère la suite $\left(c_n\right)$ définie pour tout entier naturel $n$ par : $c_n = 5u_n + 4v_n$.
		
Démontrer que la suite $\left(c_n\right)$ est constante, c'est-à-dire que pour tout entier naturel $n$, on a : $c_{n+1} = c_n$. 

En déduire que, pour tout entier naturel $n$ ,\, $c_n = 100$.
		\item Déterminer la valeur commune des limites  $\ell$ et $\ell'$.
	\end{enumerate}
\end{enumerate}

\bigskip

\end{document}
